\begin{table}[htbp]
\centering
\caption{分组异质性检验：显著组与不显著组的对比}
\label{tab:heterogeneity}
\small
\resizebox{\textwidth}{!}{%
\begin{tabular}{lccccc}
\toprule
分组方案 & 低组效应（$p$） & 高组效应（$p$） & 组间差异 & 差异$p$值 & 聚类层级 \\
\midrule
供应商关系 vs 客户关系 & -0.0033 (0.531) & 0.0295 (<0.001) & 0.0327 & <0.001 & 关系对 \\
政策前关系成熟度：其他关系 vs 长期关系 & 0.0183 (<0.001) & -0.0162 (0.309) & -0.0345 & 0.036 & 关系对 \\
RCEP前自贸协定覆盖：已有协定伙伴 vs 日本 & -0.0026 (0.605) & 0.0378 (<0.001) & 0.0403 & <0.001 & 关系对 \\
供应商关系中的企业行业：非制造业 vs 制造业 & -0.0205 (0.263) & 0.0274 (0.030) & 0.0479 & 0.031 & 关系对 \\
\bottomrule
\end{tabular}}
\begin{tablenotes}
\footnotesize 注：括号内为组内处理效应的$p$值；组间差异来自同一回归中的线性限制检验。长期关系是指截至2021年末关系年龄至少为三年的关系。企业行业分组仅使用供应商关系及政策前行业信息可匹配的企业。所有模型均含关系对和年份固定效应，标准误按关系对聚类。
\end{tablenotes}
\end{table}